% answers_short.tex —— 2004 年·普通高等学校招生考试·湖北卷·理科 · 简短答案速查
% 由 tools/make_short_answers.py 自 answers.tex 的 \gksolhead/\gkanswer 提取，
% 并与 2004/2004_数学真题_湖北理科_答案要点.md 逐题比对后生成，勿手改（改源 MD 后重跑）。

\begin{gknotice}
\textbf{使用说明}：本节只给最终答案便于快速核对，完整推导见解析版。答案为编写者独立推导，\textbf{非官方标准答案}。
\end{gknotice}

\gkpart{选择题}
\begin{center}\begin{tabular}{ccccccccccccc}
\toprule
题号 & 1 & 2 & 3 & 4 & 5 & 6 & 7 & 8 & 9 & 10 & 11 & 12 \\
\midrule
答案 & D & A & C & B & B & D & B & C & C & A & D & A \\
\bottomrule
\end{tabular}\end{center}

\gkpart{填空题}
\begin{itemize}[leftmargin=2.6em]
  \item[\textbf{13.}] $4$
  \item[\textbf{14.}] $240$
  \item[\textbf{15.}] ④
  \item[\textbf{16.}] $-1.6$（或 $-\dfrac{8}{5}$）
\end{itemize}

\gkpart{解答题}
\begin{itemize}[leftmargin=2.6em]
  \item[\textbf{17.}] $\dfrac{5\sqrt{3} - 12}{26}$
  \item[\textbf{18.}] （1）点 $F$ 为棱 $CD$ 的中点；（2）$\arccos\left(-\dfrac{1}{3}\right)$（或 $\pi - \arccos\dfrac{1}{3}$）
  \item[\textbf{19.}] 当 $\theta = 0$ 时，$\vec{BP} \cdot \vec{CQ}$ 的值最大，最大值为 $0$
  \item[\textbf{20.}] （1）$(-2, -\sqrt{2})$；（2）存在，$k = -\dfrac{6 + \sqrt{6}}{5}$
  \item[\textbf{21.}] 联合采用甲、乙两种预防措施使总费用最少，最少总费用为 $81$ 万元
  \item[\textbf{22.}] （1）$A = \dfrac{a + \sqrt{a^2 + 4}}{2}$；（2）证明见解析；（3）$a \geqslant \dfrac{3}{2}$（或 $\left[\dfrac{3}{2}, +\infty\right)$）
\end{itemize}
